\subsection*{Cap\'itulo 2: Matem\'atica discreta}
\textbf{2.1}\label{resp:2-1}\quad
(1) $4\cdot3=12$. (2) $4\cdot3\cdot2=24$. (3) $10^6=1\,000\,000$. (4) $26^3\cdot10^3=17\,576\,000$. (5) $5+3=8$. (6) $62^8=218\,340\,105\,584\,896\approx2{,}18\times10^{14}$. (7) Safari: $1\cdot2=2$; los otros 3 navegadores: $3\cdot3\cdot2=18$. Total $20$. (8) Son 25 objetos en 24 cajas: por el palomar, alguna caja recibe 2 o m\'as.

\medskip\noindent\textbf{2.2}\label{resp:2-2}\quad
(1) $4!=24$. (2) $\binom{10}4=210$. (3) $\binom{12}5=792$. (4) $P(6, 2)=30$. (5) $P(8, 3)=336$ y $\binom83=56$. (6) $\binom{10}4=210$ y $P(10, 4)=5040$. (7) $\frac{5!}{2!}=60$ (la E se repite dos veces). (8) $\binom nr=\frac{n!}{r!(n-r)!}$ es sim\'etrica en $r$ y $n-r$. $\binom{10}3=120=\binom{10}7$, la diferencia es $0$.

\medskip\noindent\textbf{2.3}\label{resp:2-3}\quad
(1) $2, 6, 18, 54, 162, 486$. (2) $4n$. (3) $2^5=32$. (4) $140$. (5) $1, 1, 2, 3, 5, 8, 13, 21$. (6) $2^5-1=31$. (7) $T(n)=\frac{n(n+1)}2$; $T(4)=10$. (8) Porque la recurrencia da el orden de crecimiento (por ejemplo $n\log n$ frente a $n^2$) sin ejecutar el algoritmo.

\medskip\noindent\textbf{2.4}\label{resp:2-4}\quad
(1) Compuesto: $7\cdot13$. (2) $2^2\cdot3^2\cdot5$. (3) $84=2\cdot30+24$, $30=1\cdot24+6$, $24=4\cdot6$: $\operatorname{MCD}=6$. (4) $36$. (5) S\'i. (6) $252=2\cdot105+42$, $105=2\cdot42+21$, $42=2\cdot21$: $21$. (7) Si $n=ab$ con $a\le b$, entonces $a\le\sqrt n$; un divisor propio siempre aparece antes de $\sqrt n$. (8) Con $n=55$ cualquiera lo factoriza al instante; RSA real usa primos de cientos de d\'igitos.

\medskip\noindent\textbf{2.5}\label{resp:2-5}\quad
(1) 5. (2) 6. (3) 3. (4) $63\bmod5=3$. (5) $20\bmod16=4$. (6) 5. (7) $13$. (8) RSA usa exponenciaci\'on modular e inversos modulares; una blockchain tambi\'en requiere hash, firmas, consenso, redes y validaci\'on.

\medskip\noindent\textbf{2.6}\label{resp:2-6}\quad
(1) Todos tienen grado 2. (2) $14$. (3) Instalar $\to$ compilar $\to$ ejecutar pruebas. (4) $\begin{bmatrix}0&1&1&0\\1&0&0&1\\1&0&0&1\\0&1&1&0\end{bmatrix}$. (5) No. (6) $A\to E\to C\to B$, costo $6$. (7) $\binom82=28$. (8) En la red f\'isica las aristas son cables reales con costo y capacidad; en las relaciones entre usuarios son v\'inculos l\'ogicos (amistad, seguimiento).

\medskip\noindent\textbf{2.7}\label{resp:2-7}\quad
(1) Una (por ejemplo $B$--$D$). (2) $\binom42=6$ cables. (3) $2^6=64$. (4) Respuesta libre; por ejemplo, un anillo de $n$ servidores donde el siguiente de $i$ es $(i+1)\bmod n$.

\medskip\noindent\textbf{Actividades de cierre:} (1) $24$. (2) $35$. (3) $\Theta(n\log n)$, como MergeSort. (4) $21$. (5) $4$. (6) a (8): respuesta libre; verifica que tu modelo identifique grados, caminos y puntos cr\'iticos.

\medskip\noindent\textbf{Autoevaluaci\'on 2:}\label{resp:q2}\quad 1 b, 2 b, 3 a, 4 b, 5 b, 6 b, 7 a, 8 c, 9 c, 10 b.
